\documentclass[11pt]{article}
\usepackage[utf8]{inputenc}
\usepackage[spanish]{babel}
\usepackage{amsmath, amssymb}
\usepackage[margin=2.7cm]{geometry}
\usepackage{booktabs}
\usepackage{hyperref}

\newcommand{\afinidad}{w}

\title{Cuantificación bidireccional de temas en proyectos y publicaciones\\
\large Fórmulas, procedencia y ejemplo trabajado}
\author{Instituto Pulso PUCP --- Libro Blanco}
\date{\today}

\begin{document}
\maketitle

\section{Planteamiento}

Cada proyecto o publicación (una \emph{unidad} $u$) trata de varios temas a la
vez, en distinto grado. El equipo pidió responder dos preguntas distintas
para el mismo par (unidad, tema):

\begin{enumerate}
  \item \textbf{¿Cuánto de la unidad es este tema?} --- si miro solo este
    proyecto, qué fracción de lo que hace corresponde a este tema.
  \item \textbf{¿Cuánto del tema aporta la unidad?} --- si miro solo este
    tema, qué fracción de todo lo que lo compone viene de este proyecto.
  \end{enumerate}

Estas dos preguntas casi nunca tienen la misma respuesta numérica, y la
diferencia depende únicamente del tamaño del tema (\S\ref{sec:ejemplo}). No
existe una etiqueta corta y estándar para este par de lecturas --- se nombran
aquí por la pregunta que responden en vez de inventar un término que se
prestaría a confusión --- pero las matemáticas de abajo sí son estándar
(\S\ref{sec:procedencia}).

\section{La matriz de afinidad}

Todo parte de una matriz de afinidad $\afinidad[u,t] \geq 0$: cuánta
evidencia hay de que la unidad $u$ trata del tema $t$. Se construye
mezclando dos fuentes, cada una normalizada a que su fila sume 1 antes de
mezclarlas:

\begin{itemize}
  \item \textbf{Extracción por LLM} (peso $\alpha = 0.7$ por defecto): un LLM
    lee el texto de la unidad y devuelve de 1 a 5 temas, cada uno con una
    cita literal verificada contra el texto y una confianza en $[0,1]$. Si
    varios temas extraídos de la misma unidad caen, al normalizar, dentro
    del mismo tema normalizado, sus confianzas se suman.
  \item \textbf{Similitud de embeddings} (peso $1-\alpha = 0.3$): coseno entre
    el embedding del texto completo de la unidad y el centroide de cada
    tema (la media de los embeddings de los temas extraídos que ese
    agrupamiento reunió). Se recorta antes de usarla --- descartar lo que no
    supera un piso de 0.30, quedarse con los $k=8$ temas más cercanos y
    elevar a una potencia $p=3$ --- porque en este corpus el coseno entre dos
    textos \emph{no} relacionados ronda 0.45, no 0: sin el recorte, toda
    unidad se repartiría casi por igual entre todos los temas y el
    coeficiente de partición (\S\ref{sec:particion}) caería a su mínimo.
\end{itemize}

Solo la primera fuente tiene evidencia textual directa detrás; la segunda es
proximidad semántica del texto completo, no algo que el LLM haya afirmado de
esa unidad en particular. Sin la segunda, con solo 2--4 temas extraídos por
unidad, la partición resultante es casi nítida en vez de difusa: cada unidad
cae en un tema casi al 100\,\% y la segunda pregunta pierde sentido (ver la
corrida \texttt{legacy-49} en \texttt{docs/cuantificacion\_temas.md}).

\section{Las dos fórmulas}

\begin{equation}
  C[u,t] \;=\; \frac{\afinidad[u,t]}{\displaystyle\sum_{t'} \afinidad[u,t']}
  \qquad\text{(``cuánto de la unidad es este tema'', } = P(t \mid u) \text{)}
\end{equation}

\noindent
Se fija la unidad $u$ y se suma $\afinidad$ sobre \emph{todos los temas}: el
total de la fila de esa unidad. Todas las $C[u,\cdot]$ de una misma unidad
suman 1.

\begin{equation}
  E[t] \;=\; \sum_{u} C[u,t]
  \label{eq:E}
\end{equation}

\noindent
El tamaño del tema en \emph{unidades-equivalentes}: se fija el tema $t$ y se
suma $C$ sobre \emph{todas las unidades} --- el total de la columna. Como
$\sum_u C[u,t] $ ya reparte 100\,\% de cada unidad entre sus temas,
$\sum_t E[t] = N$ (el total de unidades): los tamaños de los temas se pueden
sumar sin doble conteo.

\begin{equation}
  K[u,t] \;=\; \frac{C[u,t]}{E[t]}
  \qquad\text{(``cuánto del tema aporta la unidad'', } = P(u \mid t) \text{)}
\end{equation}

\noindent
Se divide la misma celda $C[u,t]$ entre el tamaño de su columna. Todas las
$K[\cdot,t]$ de un mismo tema suman 1.

\paragraph{Las dos fórmulas comparten la celda $C[u,t]$.} Lo único que
cambia es la dirección de la suma en el denominador: por fila (a través de
temas, para esta unidad) en la primera, por columna (a través de unidades,
para este tema) en la segunda. $E[t]$ es, de hecho, esa suma por columna, y
es el puente entre las dos: dividir por $E[t]$ es la regla de Bayes escrita
al revés --- la segunda lectura es la primera corregida por lo grande que sea
el tema.

\section{Complementos: cuántos cuentan de verdad}

\begin{equation}
  \text{temas efectivos de } u \;=\; \frac{1}{\sum_t C[u,t]^2}
  \qquad\qquad
  \text{unidades efectivas de } t \;=\; \frac{1}{\sum_u K[u,t]^2}
\end{equation}

\noindent
El \emph{número efectivo} $1/\sum p^2$ dice cuántas piezas, si todas
pesaran igual, producirían la misma concentración que la distribución real.
Si diez proyectos aportan 10\,\% cada uno, salen 10 efectivos; si uno aporta
90\,\% y los otros nueve 1.1\,\% cada uno, sale 1.2: dice cuántos cuentan de
verdad, no cuántos hay.

\subsection{Coeficiente de partición}
\label{sec:particion}

\begin{equation}
  PC \;=\; \frac{1}{N}\sum_{u}\sum_{t} C[u,t]^2
  \qquad\qquad
  PC_{\text{norm}} \;=\; \frac{PC - 1/T}{1 - 1/T}
\end{equation}

\noindent
$PC$ vale 1 si cada unidad cae en un único tema (partición nítida) y $1/T$
si se reparte por igual entre los $T$ temas (partición uniforme, sin
información). $PC$ crudo depende del número de temas $T$ --- un 0.45 con 49
temas y un 0.45 con 121 no significan lo mismo --- por eso se reporta
$PC_{\text{norm}}$ (normalización de Dunn/Backer), comparable entre corridas
con distinto $T$.

\begin{center}
\begin{tabular}{@{}lll@{}}
\toprule
$PC_{\text{norm}}$ & Interpretación & \\
\midrule
$> 0.95$ & partición prácticamente nítida & la contención no aporta nada sobre una asignación simple \\
$0.2$ -- $0.8$ & rango informativo & reportar \\
$< 0.05$ & prácticamente uniforme & revisar el recorte de la similitud \\
\bottomrule
\end{tabular}
\end{center}

\section{Ejemplo trabajado}
\label{sec:ejemplo}

Proyecto 767, \emph{``Stabilizability with Markov Parameters Partially
Known and Metrics of...''}, tema \emph{``Estabilidad estocástica de sistemas
singulares con saltos markovianos''}, corrida \texttt{full-proj-409}.

\paragraph{Paso 1 --- la fila de afinidades de esta unidad (las 8 con valor):}
\[
  0.8876,\; 0.0278,\; 0.0205,\; 0.0179,\; 0.0124,\; 0.0120,\; 0.0115,\; 0.0103
  \qquad \sum = 1.0000
\]

\paragraph{Paso 2 --- contención:}
\[
  C[u,t] = \frac{0.8876}{1.0000} = 0.8876 \;\to\; 88.8\,\%
\]
El 88.8\,\% de lo que trata este proyecto es este tema.

\paragraph{Paso 3 --- tamaño del tema:} sumando esa misma columna sobre las
409 unidades del dominio, el tema recibe fracciones de 17 proyectos
distintos y acumula
\[
  E[t] = 4.48 \text{ unidades-equivalentes.}
\]
Vale 4.48 proyectos enteros, aunque lo rocen 17.

\paragraph{Paso 4 --- contribución:}
\[
  K[u,t] = \frac{C[u,t]}{E[t]} = \frac{0.8876}{4.48} = 0.1982 \;\to\; 19.8\,\%
\]
Este proyecto es, él solo, el 19.8\,\% de todo el tema.

\medskip
\noindent
El mismo $\afinidad[u,t]=0.8876$ produce 88.8\,\% o 19.8\,\% según contra
qué total se mida: la fila del proyecto, o la columna del tema. Si el tema
fuera del tamaño del tema más grande de esta corrida (20.3
unidades-equivalentes, sostenido por 107 proyectos), la misma contención del
88.8\,\% se traduciría en apenas un 4.4\,\% de contribución --- por eso las
dos preguntas hacen falta por separado: preguntar ``¿cuánto pesa esta unidad
en este tema?'' sin decir en qué dirección no tiene una sola respuesta.

\section{De dónde sale esto}
\label{sec:procedencia}

\begin{itemize}
  \item \textbf{La matriz de pertenencia difusa.} Que una unidad pertenezca
    a varios temas a la vez, en grados que suman 1, es una \emph{partición
    difusa}: el objeto central del agrupamiento difuso (\emph{fuzzy
    c-means}, Bezdek, 1981). De ahí viene también el coeficiente de
    partición que sugirió el profesor, que no es el número por celda sino
    un índice de validez del conjunto completo.
  \item \textbf{Las dos direcciones.} Son probabilidad condicionada
    elemental sobre una distribución conjunta $P(u,t) \propto \afinidad[u,t]$,
    unidas por la regla de Bayes. $P(t\mid u)$ es, además, lo mismo que
    devuelve un modelo de tópicos tipo LDA como distribución
    tema-documento; la novedad aquí no es esa mitad, sino publicar también
    la otra.
  \item \textbf{Las unidades-equivalentes.} Son \emph{conteo fraccionado}
    (\emph{fractional counting}), la práctica estándar en bibliometría para
    repartir el crédito de una publicación entre sus autores o
    instituciones en vez de contarla entera para cada uno. Aquí se reparte
    entre temas en vez de entre autores.
  \item \textbf{El número efectivo} $1/\sum p^2$ se reinventa en todas
    partes con el mismo nombre distinto: inverso del índice de Simpson en
    ecología, inverso del índice de Herfindahl--Hirschman en economía de la
    competencia, ``número efectivo de partidos'' en ciencia política.
\end{itemize}

\section{Reproducibilidad}

Todas las cifras de este documento salen de la corrida \texttt{full-proj-409}
del dominio \texttt{projects} (409 proyectos, filtrados por universo mínimo
declarado: cerrado y 2010+, con \texttt{title}, \texttt{knowledge\_area} y al
menos una línea de investigación). Fórmulas implementadas en
\texttt{scripts/lib/lb\_membership.py}; corren sobre Amazon Bedrock
(\texttt{Claude Sonnet 5} para extracción y nombrado,
\texttt{cohere.embed-multilingual-v3} para embeddings) y se guardan en
PostgreSQL con \texttt{pgvector} en Amazon RDS. El explorador interactivo de
esta misma corrida está en el atlas temático publicado por el equipo.

\end{document}
